\documentclass[10pt,a4paper]{amsart}
\usepackage[margin=19mm]{geometry}
\usepackage{amsmath,amssymb,mathtools}
\usepackage[scr=rsfs,cal=euler]{mathalfa}
\usepackage{xfrac}
\usepackage[unicode,hidelinks,bookmarks=false]{hyperref}
\newcommand{\ie}{i.e.\ }
\newcommand{\cf}{cf.\ }
\newcommand{\resp}{respectively\ }
\newcommand{\Subsection}{\subsection}
\providecommand{\nobreakdash}{\nobreak\hbox{-}}
\setlength{\emergencystretch}{2em}
\multlinegap0pt
\usepackage{amscd,mathtools}
\usepackage{eucal}
\usepackage{xy}
\usepackage{graphpap,color}
\usepackage{cancel}
\usepackage{verbatim}
\usepackage{adjustbox}
\usepackage{pgfplots}
\pgfplotsset{compat=1.18}
\usepackage{tikz}
\usepackage{float}
\usepackage{booktabs}
\usepackage{multirow}
\usepackage{tabularx}
\usepackage{enumitem}
\setlength{\evensidemargin}{\oddsidemargin}
\numberwithin{equation}{section}
\newtheorem{theo}{Theorem}[section]
\newtheorem{coro}[theo]{Corollary}
\newtheorem{prop}[theo]{Proposition}
\newtheorem{lemm}[theo]{Lemma}
\newtheorem{conj}[theo]{Conjecture}
\newtheorem{assu}[theo]{Assumption}
\newtheorem{defi}[theo]{Definition}
\newtheorem{ques}[theo]{Question}
\newtheorem{exam}[theo]{Example}
\newtheorem{rema}[theo]{Remark}
\newtheorem{prob}[theo]{Problem}
\newtheorem{solu}[theo]{Solution}
\def\blue{\textcolor{blue}}
\def\green{\textcolor{green}}
\def\red{\textcolor{red}}
\def\black{\textcolor{black}}
\def\cyan{\textcolor{cyan}}
\newcommand{\A }{\mathbb{A}}
\newcommand{\CC}{\mathbb{C}}
\newcommand{\EE}{\mathbb{E}}
\newcommand{\LL}{\mathbb{L}}
\newcommand{\PP}{\mathbb{P}}
\newcommand{\QQ}{\mathbb{Q}}
\newcommand{\RR}{\mathbb{R}}
\newcommand{\ZZ}{\mathbb{Z}}
\newcommand{\NN}{\mathbb{N}}
\newcommand{\HH}{\mathbb{H}}
\newcommand{\bL}{\mathbf{L}}
\newcommand{\bT}{\mathbf{T}}
\newcommand{\bk}{\mathbf{k}}
\newcommand{\bm}{\mathbf{m}}
\newcommand{\bn}{\mathbf{n}}
\newcommand{\bp}{\mathbf{p}}
\newcommand{\bq}{\mathbf{q}}
\newcommand{\bs}{\mathbf{s}}
\newcommand{\bt}{\mathbf{t}}
\newcommand{\bw}{\mathbf{w}}
\newcommand{\bx}{\mathbf{x}}
\newcommand{\kk}{\bk}
\def\bA{\mathbf{A}}
\newcommand{\bP}{\mathbf{P}}
\newcommand{\bX}{\mathbf{X}}
\newcommand{\bS}{\mathbf{S}}
\newcommand{\bM}{\mathbf{M}}
\newcommand{\bH}{\mathbf{H}}
\newcommand{\bC}{\mathbf{C}}
\def\cA{{\mathcal A}}
\def\cB{{\mathcal B}}
\def\cC{{\mathcal C}}
\def\cD{{\mathcal D}}
\def\cE{{\mathcal E}}
\def\cF{{\mathcal F}}
\def\cG{{\mathcal G}}
\def\cH{{\mathcal H}}
\def\cK{{\mathcal K}}
\def\cL{{\mathcal L}}
\def\cM{{\mathcal M}}
\def\cN{{\mathcal N}}
\def\cO{{\mathcal O}}
\def\cP{{\mathcal P}}
\def\cQ{{\mathcal Q}}
\def\cR{{\mathcal R}}
\def\cS{{\mathcal S}}
\def\cT{{\mathcal T}}
\def\cU{{\mathcal U}}
\def\cV{{\mathcal V}}
\def\cW{{\mathcal W}}
\def\cX{{\mathcal X}}
\def\cY{{\mathcal Y}}
\def\cZ{{\mathcal Z}}
\def\cI{{\mathcal I}}
\def\cJ{{\mathcal J}}
\def\fB{\mathfrak{B}}
\def\fC{\mathfrak{C}}
\def\fD{\mathfrak{D}}
\def\fE{\mathfrak{E}}
\def\fF{\mathfrak{F}}
\def\fM{\mathfrak{M}}
\def\fV{\mathfrak{V}}
\def\fX{\mathfrak{X}}
\def\fY{\mathfrak{Y}}
\def\ff{\mathfrak{f}}
\def\fm{\mathfrak{m}}
\def\fp{\mathfrak{p}}
\def\ft{\mathfrak{t}}
\def\fu{\mathfrac{u}}
\def\fv{\mathfrak{v}}
\def\frev{\mathfrak{rev}}
\newcommand{\tC}{\tilde{C}}
\newcommand{\tD}{\tilde{D}}
\newcommand{\tE}{\tilde{E}}
\newcommand{\tF}{\tilde{F}}
\newcommand{\tK}{\tilde{K}}
\newcommand{\tL}{\tilde{L}}
\newcommand{\tT}{\tilde{T}}
\newcommand{\tY}{\tilde{Y}}
\newcommand{\hD}{\hat{D}}
\newcommand{\hL}{\hat{L}}
\newcommand{\hT}{\hat{T}}
\newcommand{\hX}{\hat{X}}
\newcommand{\hY}{\hat{Y}}
\newcommand{\hZ}{\hat{Z}}
\newcommand{\hf}{\hat{f}}
\newcommand{\hu}{\hat{u}}
\newcommand{\hDe}{\hat{Delta}}
\def\begeq{\begin{equation}}
\def\endeq{\end{equation}}
\def\and{\quad{\rm and}\quad}
\def\lra{\longrightarrow }
\def\mapright#1{\,\smash{\mathop{\lra}\limits^{#1}}\,}
\def\mapleft#1{\,\smash{\mathop{\longleftarrow}\limits^{#1}}\,}
\def\sO{{\mathscr O}}
\def\beq{\begin{equation}}
\def\eeq{\end{equation}}
\def\DM{Deligne-Mumford }
\def\wt{\mathrm{wt}}
\def\GG{\mathbb{G} }
\def\Gm{\mathbb{G}_m}
\def\Pic{\mathrm{Pic} }
\def\spec{\mathrm{Spec}\, }
\def\Map{\mathrm{Map} }
\def\oM{\overline{M} }
\def\Hom{\mathrm{Hom} }
\def\bzero{\mathbf{0}}
\def\bone{\mathbf{1}}
\def\bcM{\overline{\cM}}
\def\cAn{\cA^{(n)}}
\def\Yn{Y^{(n)}}
\def\bHn{\bH^{(n)}}
\setcounter{tocdepth}{1}
\begin{document}
\title[Hodge-Deligne and Poincar'e polynomials of overlinemathcal M]{Hodge-Deligne and Poincar\'e polynomials of $\overline{\mathcal M}_{1,n}$}
\author{Young-Hoon Kiem}
\date{}
\maketitle
\noindent\emph{Source and licence.} Hodge-Deligne and Poincar\'e polynomials of $\overline{\mathcal M}_{1,n}$. Version arXiv:2609.12507v1 (CC BY 4.0). This material is distributed under the Creative Commons Attribution 4.0 International licence, http://creativecommons.org/licenses/by/4.0/, as recorded in the arXiv OAI licence metadata for this exact version and shown on the abstract page https://arxiv.org/abs/2609.12507v1. Full rights notice: licenses/arxiv-2609.12507-CC-BY-4.0.txt.\par\smallskip
\noindent\textit{Editorial scope.} Counted unit: the original \texttt{theo} environment \texttt{2} at source lines 338--350 together with its complete proof at lines 351--353; the delivered document keeps source lines 322--354 so that every referenced label and every citation resolves inside it. Native editable source is the paper's own concatenated TeX; the AMS wrapper is editorial formatting only. No publisher class, upstream script or unrelated artwork is redistributed.
\section{Moduli spaces of weighted pointed curves and Hodge-Deligne polynomials}
In this section, we collect necessary facts about Hassett moduli spaces \cite{Hst} and Hodge-Deligne polynomials that will be used throughout this paper. 

\begin{defi} \cite{Hst}
Let $g, n\ge 0$. For $\cA=(a_1,\cdots,a_n)\in (0,1]^n$ with $2g-2+\sum_i a_i>0$, an integral curve $C$ with marked points $p_1, \cdots,p_n\in C$ is called $\cA$-\emph{stable} if \begin{enumerate}
\item $C$ has at worst nodal singularities and $p_i$ are all smooth points;
\item if $p_{i_1}=p_{i_2}=\cdots=p_{i_k}$, then $a_{i_1}+\cdots+a_{i_k}\le 1$;
\item $\omega_C(a_1p_1+\cdots +a_np_n)$ is ample. 
\end{enumerate}
\end{defi}
Let $\bcM_{g,\cA}$ denote the moduli stack of $\cA$-stable $n$-pointed curves of genus $g$. 
In particular, when $a_i>1/2$ for all $i$, $\bcM_{g,\cA}$ coincides with the Knudsen-Deligne-Mumford moduli stack $\bcM_{g,n}$ of stable curves of genus $g$ with $n$ marked points \cite{DM,Knu}.

\medskip

Let us summarize useful facts about $\bcM_{g,\cA}$ from \cite{Hst}.

\begin{theo}\label{2} 
Let $\cA=(a_1,\cdots,a_n)\in (0,1]^n$ with $2g-2+\sum_{i=1}^n a_i>0$.

(1) $\bcM_{g,\cA}$ is a smooth connected proper Deligne-Mumford stack of dimension $3g-3+n$ with projective coarse moduli space.  

(2) If $\cA=(a_1,\cdots,a_n)$, $\cA'=(a'_1,\cdots, a'_n)$ and $a_i\ge a'_i$ for all $i$, then we have a natural morphism 
\beq\label{1} \rho_{\cA,\cA'}:\bcM_{g,\cA}\lra \bcM_{g,\cA'}\eeq 
which is obtained by contracting rational components in an $\cA$-stable curve which are not $\cA'$-stable. 

(3) For $2<r\le n$, if $a'_{i_1}+\cdots+a'_{i_r}\le 1$  and $a_{i_1}+\cdots+a_{i_r}>1$ but any proper subset  of $\{a_{i_1},\cdots, a_{i_r}\}$ has sum at most one while $a_j=a_j'$ for $j\ne i_1,\cdots,i_r$, then the morphism $\rho_{\cA,\cA'}$ is the smooth blowup along the locus $\bcM_{g,\bar \cA'}$ in $\bcM_{g,\cA'}$ where $p_{i_1}=\cdots=p_{i_r}$ and $\bar \cA'$ is obtained from $\cA'$ by replacing $a'_{i_1}, \cdots, a'_{i_r}$ by their sum. 

(4) If $r\le 2$ for any $I\subset [n]$ such that $a_{i_1}+\cdots+a_{i_r}>1$ and $a'_{i_1}+\cdots+a'_{i_r}\le 1$, then $\rho_{\cA,\cA'}$ is an isomorphism. 
\end{theo}

\begin{proof}
    (1) is \cite[Theorem 2.1]{Hst} and (2) is \cite[Theorem 4.1]{Hst} while (3) is \cite[Remark 4.6]{Hst} and (4) is \cite[Corollary 4.7]{Hst}. 
\end{proof}

\begin{thebibliography}{99}
\bibitem[DM]{DM} P. Deligne and D. Mumford. {\em The irreducibility of the space of curves of given genus}. Publ. Math. Inst. Hautes Études Sci. (36) (1969) 75–109.
\bibitem[Hst]{Hst} B. Hassett. {\em Moduli spaces of weighted pointed stable curves.} Adv. Math. 173 (2003) 316–352.
\bibitem[Knu]{Knu} F. Knudsen. {\em The projectivity of the moduli space of stable curves, II}. Math. Scand. \textbf{52} (1983), 161--199.
\end{thebibliography}
\end{document}
